\documentclass[10pt,a4paper]{article}
\usepackage[utf8]{inputenc}
\usepackage[spanish,es-noquot]{babel}
\usepackage{amsmath,amssymb}
\usepackage[margin=1cm]{geometry}
\usepackage{xcolor}
\usepackage{multicol}
\usepackage{enumitem}

\spanishdeactivate{<>}
\pagestyle{empty}

% Definición de colores
\definecolor{mainblue}{rgb}{0.1, 0.4, 0.7}
\definecolor{boxbg}{rgb}{0.93, 0.95, 0.98}

% Estilo de cajas de sección
\newcommand{\sectionbox}[1]{%
  \vspace{0.15cm}\noindent
  \fcolorbox{mainblue}{boxbg}{\parbox{\dimexpr\linewidth-2\fboxsep-2\fboxrule\relax}{\textbf{\color{mainblue}\small #1}}}%
  \vspace{0.1cm}%
}

\setlist[itemize]{nosep, leftmargin=1.2em}
\setlist[enumerate]{nosep, leftmargin=1.2em}

\begin{document}

\begin{center}
    {\Large \textbf{\color{mainblue} RESUMEN DE EXAMEN: POLINOMIOS Y FRACCIONES ALGEBRAICAS}}\\[0.1cm]
    {\small \textbf{Ficha de Repaso Rápido --- Fórmulas, Teoremas y Métodos}}
\end{center}

\begin{multicols*}{2}

% --------------------------------------------------
\sectionbox{1. VALOR NUMÉRICO E IDENTIDADES NOTABLES}
% --------------------------------------------------
\begin{itemize}
    \item \textbf{Valor Numérico $P(a)$:} Sustituir la variable $x$ por el número $a$. Si $P(a) = 0$, $a$ es una \textbf{raíz} de $P(x)$.
    \item \textbf{Identidades Notables Clave:}
    \begin{align*}
        (a + b)^2 &= a^2 + 2ab + b^2 \\
        (a - b)^2 &= a^2 - 2ab + b^2 \\
        (a + b)(a - b) &= a^2 - b^2
    \end{align*}
\end{itemize}

% --------------------------------------------------
\sectionbox{2. DIVISIÓN DE POLINOMIOS Y RUFFINI}
% --------------------------------------------------
\begin{itemize}
    \item \textbf{Algoritmo de la División:} $D(x) = d(x) \cdot C(x) + R(x)$ donde $\text{grado}(R) < \text{grado}(d)$.
    \item \textbf{Regla de Ruffini:} Se usa cuando el divisor es del tipo $(x - a)$ o $(x + a)$.
    \begin{itemize}
        \item Completar los coeficientes faltantes con \textbf{ceros}.
        \item Si se divide entre $(x + a)$, en Ruffini se usa el número \textbf{$-a$}.
    \end{itemize}
    \item \textbf{Candidatos a raíces enteras:} Divisores del término independiente de $P(x)$.
\end{itemize}

% --------------------------------------------------
\sectionbox{3. TEOREMA DEL RESTO Y DEL FACTOR}
% --------------------------------------------------
\begin{itemize}
    \item \textbf{Teorema del Resto:} El resto de la división $P(x) : (x - a)$ coincide con el valor numérico $P(a)$.
    \[
    \text{Resto } R = P(a)
    \]
    \item \textbf{Teorema del Factor:} $P(x)$ es divisible por $(x - a) \iff P(a) = 0$.
    \item \textbf{Cálculo de parámetros ($m$):} Igualar $P(a)$ al resto dado (o a $0$ si es divisible) y despejar $m$.
\end{itemize}

% --------------------------------------------------
\sectionbox{4. GUÍA PASO A PASO PARA FACTORIZAR}
% --------------------------------------------------
Para descomponer en factores un polinomio $P(x)$:
\begin{enumerate}
    \item \textbf{Sacar factor común:} Extraer $x$ (con el menor exponente) y/o factores numéricos.
    \[
    3x^3 - 12x^2 + 9x = 3x(x^2 - 4x + 3)
    \]
    \item \textbf{Identidades Notables:} Identificar cuadrados perfectos o diferencia de cuadrados.
    \[
    x^2 - 9 = (x - 3)(x + 3), \quad x^2 + 4x + 4 = (x + 2)^2
    \]
    \item \textbf{Regla de Ruffini:} Probar divisores del término independiente hasta obtener resto $0$.
    \item \textbf{Ecuación de 2.º Grado:} Para trinomios $ax^2 + bx + c$, resolver $ax^2 + bx + c = 0 \implies a(x - x_1)(x - x_2)$.
\end{enumerate}

\columnbreak

% --------------------------------------------------
\sectionbox{5. FRACCIONES ALGEBRAICAS}
% --------------------------------------------------
\begin{itemize}
    \item \textbf{Simplificación:} Factorizar numerador y denominador y simplificar los factores comunes repetidos.
    \[
    \frac{x^2 - 1}{x^2 - 2x + 1} = \frac{(x-1)(x+1)}{(x-1)^2} = \frac{x + 1}{x - 1}
    \]
    \item \textbf{Suma y Resta:} Hallar el $\text{m.c.m.}$ de los denominadores tras haberlos factorizado.
    \item \textbf{Multiplicación:} 
    \[
    \frac{A}{B} \cdot \frac{C}{D} = \frac{A \cdot C}{B \cdot D} \quad \text{(\small\textbf{¡Factorizar antes de multiplicar!})}
    \]
    \item \textbf{División (producto en cruz):} 
    \[
    \frac{A}{B} : \frac{C}{D} = \frac{A \cdot D}{B \cdot C}
    \]
\end{itemize}

% --------------------------------------------------
\sectionbox{6. TRADUCCIÓN ALGEBRAICA Y GEOMETRÍA}
% --------------------------------------------------
\begin{itemize}
    \item \textbf{Expresiones Frecuentes:}
    \begin{itemize}
        \item Mitad de $y$ más 6: $\dfrac{y}{2} + 6$
        \item Inversos de dos números: $\dfrac{1}{x} - \dfrac{1}{y}$
        \item Números consecutivos: $x, \quad x + 1$
        \item Algoritmo división: $D = d \cdot c + r$
    \end{itemize}
    \item \textbf{Fórmulas Geométricas:}
    \begin{itemize}
        \item \textbf{Teorema de Pitágoras:} $h = \sqrt{c_1^2 + c_2^2} \implies c_2 = \sqrt{h^2 - c_1^2}$
        \item \textbf{Rectángulo:} $P = 2b + 2h$, \quad $A = b \cdot h$
        \item \textbf{Triángulo Rectángulo:} $A = \dfrac{\text{cateto}_1 \cdot \text{cateto}_2}{2}$
        \item \textbf{Área No Edificada / Ajardinada:}
        \[
        A_{\text{libre}} = A_{\text{total}} - A_{\text{edificada}}
        \]
    \end{itemize}
\end{itemize}

% --------------------------------------------------
\sectionbox{7. ¡ERRORES TÍPICOS A EVITAR!}
% --------------------------------------------------
\begin{itemize}
    \item $\mathbf{(a + b)^2 \neq a^2 + b^2}$ $\longrightarrow$ ¡Falta el doble producto $2ab$!
    \item $\mathbf{-(x - 3) \neq -x - 3}$ $\longrightarrow$ El signo negativo cambia el signo de \textbf{todo} el paréntesis: $-x + 3$.
    \item No simplificar términos en sumas: $\dfrac{x + 2}{x} \neq 2$. Solo se simplifican factores multiplicando.
\end{itemize}

\end{multicols*}

\end{document}